
\begin{table}[H]
\RawCaption%
  {\caption[]{Summary statistics of non-marginalized and marginalized municipalities, Mexico (\%)}
   \label{t:descriptives}}
\floatbox[{\captop}]{table}[\FBwidth]
  {}
  {\scalebox{0.9}{
      \begin{tabular}{lcc}
\hline
 & \textbf{Non-marginalized areas} & \textbf{Marginalized areas} \\
\hline

\multicolumn{3}{l}{\textbf{A. Progresa program intensity}} \\
\quad 1999 & 0.10 (0.13) & 0.44 (0.22) \\
\quad 2005 & 0.28 (0.19) & 0.72 (0.17) \\[6pt]

\multicolumn{3}{l}{\textbf{B. Characteristics of municipalities, 1990}} \\
\multicolumn{3}{l}{\quad \textit{Components of the marginality index}} \\
\quad Illiterate population & 13.33 (6.57) & 33.33 (13.3) \\
\quad With less than primary education & 45.74 (12.32) & 69.56 (9.68) \\
\quad Without a toilet & 25.94 (16.62) & 60.30 (18.54) \\
\quad Without electricity & 11.81 (10.29) & 36.39 (24.72) \\
\quad Without running water & 19.43 (14.87) & 50.65 (24.07) \\
\quad With crowding & 59.84 (9.79) & 73.96 (8.24) \\
\quad With dirt floor & 22.05 (14.21) & 61.98 (21.58) \\
\quad In communities with $<$5{,}000 inhabitants & 60.18 (36.15) & 95.50 (13.54) \\
\quad Earning $<$2× minimum wage & 69.19 (11.60) & 85.82 (8.55) \\[6pt]

Number of municipalities & 1,234 & 1,119 \\
\hline
\end{tabular}

    }
    \vspace{-3mm}
    \floatfoot{Notes: Section A reports the mean proportion (\%) with standard deviations of municipalities with a ratio of the Progresa beneficiary households to the total number of households in the municipality. Section B presents sample mean (\%) of each variable and standard deviations are in parentheses. Marginalized areas include municipalities categorized as high and very high level of the Margination Index, while non-marginalized areas include municipalities categorized as medium, low, and very low level of the Margination Index. “With crowding” is measured by number of rooms divided by household size. Source: Authors’ calculations using Mexican Census 1990 and Progresa administrative data on beneficiaries.}
  }
\end{table}


\begin{table}[H]
\RawCaption%
  {\caption[]{Effects of Progresa on Elderly Mortality (65+) by Sex, Marginalized Municipalities}
   \label{t:did}}
\floatbox[{\captop}]{table}[\FBwidth]
  {}
  {\scalebox{0.9}{
      \begin{tabular}{lcccc} \hline \hline
& \multicolumn{1}{c}{(1)} & \multicolumn{1}{c}{(2)} & \multicolumn{1}{c}{(3)} & \multicolumn{1}{c}{(4)} \\ \toprule
\underline{\textit{Panel A: Pooled}}  \\ 
\textit{Intensity 1999 x post (1997-2006)} &        0.369 &        0.354 &       -0.777 &       -0.719 \\ 
 & (       2.546) & (       2.547) & (       1.464) & (       1.464) \\ 
  & & & & \\ 
\textit{Intensity 2005 x post (1997-2006)} &       -5.525* &       -5.513* &       -2.787 &       -2.796 \\ 
 & (       3.075) & (       3.074) & (       1.806) & (       1.803) \\ 
  & & & & \\ 
Mean (1991-1996) &        46.47 &        46.47 &        46.47 &        46.47 \\ 
Obs &       17,904 &       17,904 &       17,904 &       17,904 \\ 
  & & & & \\ 
\underline{\textit{Panel B: Males}}  \\ 
\textit{Intensity 1999 x post (1997-2006)} &       -1.627 &       -1.588 &       -2.288 &       -2.206 \\ 
 & (       2.767) & (       2.776) & (       1.732) & (       1.735) \\ 
  & & & & \\ 
\textit{Intensity 2005 x post (1997-2006)} &       -1.488 &       -1.519 &        0.214 &        0.201 \\ 
 & (       3.309) & (       3.312) & (       2.000) & (       1.994) \\ 
  & & & & \\ 
Mean (1991-1996) &        47.26 &        47.26 &        47.26 &        47.26 \\ 
Obs &       17,904 &       17,904 &       17,904 &       17,904 \\ 
  & & & & \\ 
\underline{\textit{Panel C: Females}}  \\ 
\textit{Intensity 1999 x post (1997-2006)} &        1.737 &        1.668 &        0.592 &        0.620 \\ 
 & (       2.932) & (       2.928) & (       1.824) & (       1.824) \\ 
  & & & & \\ 
\textit{Intensity 2005 x post (1997-2006)} &       -9.117** &       -9.063** &       -5.637** &       -5.641** \\ 
 & (       3.776) & (       3.770) & (       2.205) & (       2.205) \\ 
  & & & & \\ 
Mean (1991-1996) &        46.11 &        46.11 &        46.11 &        46.11 \\ 
Obs &       17,904 &       17,904 &       17,904 &       17,904 \\ 
  & & & & \\ 
No. Mun &        1,119 &        1,119 &        1,119 &        1,119 \\ 
  & & & & \\ 
Year FE & Y & Y & Y & Y \\ 
Mun FE & Y & Y & Y & Y \\ 
Seguro Popular & N & Y & N & Y \\ 
Weights & N & N & Y & Y \\ 
Cluster SE: Mun & Y & Y & Y & Y \\ 
\bottomrule
\end{tabular}
    }
    \vspace{-3mm}
    \floatfoot{Notes: The analysis is restricted to marginalized municipalities, 1991–2006. Progresa and Seguro Popular intensities are continuous variables. The post-dummy equals 1 for years 1997–2006. Standard errors are clustered at the municipality level and shown in parentheses. ***p$<0.01$, **p$<0.05$, *p$<0.1$.
}
  }
\end{table}


\begin{table}[htbp]
\centering
\caption{Progresa Impacts on Labor Supply and Living Arrangements: Eligible Elderly (Age 65+ in 1997)}
\label{tab:combined_elderly}
\begin{threeparttable}

\begin{tabular}{lcccc}
\toprule
 & (1) & (2) & (3) & (4) \\
 & Weekly Hours & Live Alone & With Children & Only Elderly \\
\midrule
\multicolumn{5}{l}{\textbf{Panel A: Females}} \\
\midrule
Treat $\times$ 1998 & 0.124 & 0.000 & 0.001 & $-$0.008 \\
                    & (1.136) & (0.008) & (0.011) & (0.008) \\[4pt]
Treat $\times$ 1999 & 0.292 & 0.007 & 0.015 & 0.012 \\
                    & (1.070) & (0.010) & (0.018) & (0.014) \\[4pt]
Observations        & 3,338 & 3,597 & 3,597 & 3,597 \\
Control Mean (1997) & 3.923 & 0.044 & 0.708 & 0.145 \\
\midrule
\multicolumn{5}{l}{\textbf{Panel B: Males}} \\
\midrule
Treat $\times$ 1998 & $-$1.210 & 0.012 & $-$0.008 & $-$0.007 \\
                    & (1.895) & (0.008) & (0.011) & (0.010) \\[4pt]
Treat $\times$ 1999 & 0.378 & 0.008 & 0.040$^{**}$ & $-$0.001 \\
                    & (1.872) & (0.012) & (0.020) & (0.017) \\[4pt]
Observations        & 3,170 & 3,419 & 3,419 & 3,419 \\
Control Mean (1997) & 25.288 & 0.049 & 0.676 & 0.159 \\
\bottomrule
\end{tabular}

\begin{tablenotes}
\small
\item \textit{Notes:} Difference-in-differences estimates with 1997 as baseline. Sample restricted to 
eligible households. Columns (1) reports weekly hours worked. Columns (2)--(4) report living 
arrangement outcomes: living alone, living with children, and living only with other elderly, 
respectively. Standard errors in parentheses.
\item $^{*}$ $p < 0.10$, $^{**}$ $p < 0.05$, $^{***}$ $p < 0.01$
\end{tablenotes}
\end{threeparttable}

\end{table}